\documentclass[11pt]{beamer}
\usepackage[utf8]{inputenc}
\usepackage[T1]{fontenc}
\usepackage[french]{babel}
\usepackage{graphicx}
\usepackage{xcolor}
\usepackage{hyperref}

\usetheme{Madrid}
\usecolortheme{default}

\title[Visualisations]{Stack technique des visualisations frontend}
\subtitle{Projet Full-Stack - Analyse de Performance Sportive avec IA}
\institute[Projet Informatique]{Projet Informatique}
\date{\today}

\begin{document}

\begin{frame}{Stack technique des visualisations}
\begin{table}
\centering
\scriptsize
\begin{tabular}{|l|l|l|}
\hline
\textbf{Visualisation} & \textbf{Techno / Librairie} & \textbf{Fichier} \\
\hline
Trajectoire du point & React + SVG inline + \texttt{<animate>} & \texttt{PointTrajectory} \\
\hline
Sketch Search & React + SVG inline & \texttt{sketch/page} \\
\hline
Momentum du match & React + SVG (paths, gradients) & \texttt{MomentumChart} \\
\hline
Profil serveur & React + SVG (jauges, heatmap) & \texttt{ServerProfilePanel} \\
\hline
Face-à-face (radar) & React + SVG (polygones) & \texttt{compare/page} \\
\hline
Carte embeddings 2D & \textcolor{blue}{Recharts} (ScatterChart) & \texttt{EmbeddingChart} \\
\hline
Carte embeddings 3D & \textcolor{blue}{Three.js} + react-three/fiber & \texttt{EmbeddingChart3D} \\
\hline
Lecteur vidéo & HTML5 \texttt{<video>} + SVG & \texttt{VideoPlayer} \\
\hline
\end{tabular}
\end{table}

\begin{alertblock}{Philosophie}
\small
\textbf{5 visus sur 8 en SVG inline pur} (zéro dépendance charting).\\
Seules exceptions : Recharts (scatter 2D) et Three.js (3D WebGL).
\end{alertblock}
\end{frame}

\end{document}
